\chapter{Kommentierte Quellen}

Achtundzwanzig Werke, eingeordnet: Zu jedem steht, wofür man es aufschlägt und was es leistet, das andere nicht leisten.

\lit{1}{Serre, J.-P.: \textit{Linear Representations of Finite Groups}. Springer 1977 (frz.\ 1967).}{Die Teile I--III dieses Bandes auf siebzig Seiten, in der Dichte, die Serre eigen ist: Charaktere, Orthogonalität, Induktion, Mackey, Brauers Satz --- geschrieben für Chemiker in Paris, die er unterrichtete. \textbf{Die Fundstelle für die Auslassungen~1, 2, 3 und~8}; der zweite und dritte Teil (rationale Darstellungen, Brauer-Theorie) sind das, was dieser Band nicht getan hat.}

\lit{2}{Fulton, W., Harris, J.: \textit{Representation Theory: A First Course}. Springer 1991.}{Endliche Gruppen in hundert Seiten, dann Lie-Algebren und -Gruppen in fünfhundert, immer an Beispielen: $S_n$, $SL_2$, $SL_3$, dann die Klassifikation. \textbf{Das Buch für den Weg von Kapitel~14 zu den Lie-Algebren}; die Fundstelle für Auslassungen~2, 7, 14 und~19. Die Übungen enthalten die Hälfte des Stoffs.}

\lit{3}{James, G., Liebeck, M.: \textit{Representations and Characters of Groups}. 2.\ Aufl., Cambridge 2001.}{Das langsamste und gründlichste Lehrbuch der Charaktertheorie: jede Tafel bis zur Ordnung $32$ berechnet, jede Technik an Beispielen. Ein eigenes Kapitel über Molekülschwingungen. \textbf{Für Leser, die nach Kapitel~5 noch zwanzig Tafeln selbst bauen wollen.}}

\lit{4}{Isaacs, I.\ M.: \textit{Character Theory of Finite Groups}. Academic Press 1976; Dover 1994.}{Die Charaktertheorie als Werkzeug der Gruppentheorie: Burnsides $p^{a}q^{b}$-Satz, Frobenius-Gruppen, Brauers Sätze, Clifford-Theorie, Anfänge der modularen Theorie. \textbf{Schwerer als Serre, vollständiger}; die Fundstelle für Auslassungen~1 und~3 und für Brauers Probleme aus Kapitel~4.}

\lit{5}{Etingof, P., u.\,a.: \textit{Introduction to Representation Theory}. AMS 2011.}{Aus einem MIT-Kurs für Studienanfänger: Darstellungen von Algebren, Köcher, dann endliche Gruppen, $S_n$ und Schur--Weyl. \textbf{Die Fundstelle für den Doppelkommutantensatz (Auslassung~14)} in einer Fassung ohne Vorkenntnisse; frei verfügbar.}

\lit{6}{Sagan, B.: \textit{The Symmetric Group}. 2.\ Aufl., Springer 2001.}{Kapitel~8 in voller Länge: Specht-Moduln, Verzweigung, Hakenformel mit dem Hakenspaziergang, Robinson--Schensted, Murnaghan--Nakayama, symmetrische Funktionen. \textbf{Die Fundstelle für die Auslassungen~4 bis~7}; kombinatorisch, mit Bildern, ohne Umwege.}

\lit{7}{James, G., Kerber, A.: \textit{The Representation Theory of the Symmetric Group}. Addison-Wesley 1981.}{Das Nachschlagewerk zu $S_n$, mit Tafeln bis $S_{10}$ und dem, was Sagan weglässt: modulare Theorie, Verallgemeinerungen, Geschichte. \textbf{Dort steht James' Vermutung in ihrer ursprünglichen Form} (Kapitel~8, Bruchfall).}

\lit{8}{Diaconis, P.: \textit{Group Representations in Probability and Statistics}. IMS Lecture Notes 1988.}{Fourier-Analyse auf $S_n$ für Rangdaten und Mischvorgänge --- Kapitel~12 als Statistik, mit den Wahldaten der American Psychological Association und den Riffelmischungen. \textbf{Das Buch, das zeigt, dass Charaktertafeln Daten auswerten}; frei verfügbar.}

\lit{9}{Bröcker, T., tom Dieck, T.: \textit{Representations of Compact Lie Groups}. Springer 1985.}{Teil~IV vollständig und streng: Haar-Maß, Peter--Weyl, Weyls Integralformel, maximale Tori, Charakterformel, $SU(2)$ und $SO(3)$ als Beispiele. \textbf{Die Fundstelle für die Auslassungen~8 bis~11 und~13}; zwei Göttinger Autoren, englisch geschrieben.}

\lit{10}{Hall, B.\ C.: \textit{Lie Groups, Lie Algebras, and Representations}. 2.\ Aufl., Springer 2015.}{Matrixgruppen ohne Mannigfaltigkeiten: Exponentialabbildung, Lie-Algebra, $SU(2)$, $SU(3)$, dann die allgemeine Theorie. \textbf{Der leichteste Einstieg in das, was Kapitel~11 mit Leiteroperatoren nur angedeutet hat}; die Fundstelle für Auslassung~12.}

\lit{11}{\sloppy Hein, W.: \textit{Einführung in die Struktur- und Darstellungs\-theorie der klassischen Gruppen}. Springer, Berlin 1990.}{\textbf{Auf Deutsch}: $GL_n$, $O_n$, $Sp_n$, $U_n$ und ihre Darstellungen über Tensoren, Young-Symmetrisierer und Schur--Weyl --- der Weg von Kapitel~12 zu den klassischen Gruppen in der Sprache Weyls. Knapp, mit allen Beweisen.}

\lit{12}{Böhm, M.: \textit{Lie-Gruppen und Lie-Algebren in der Physik}. Springer 2011.}{\textbf{Auf Deutsch, für Physiker}: Drehgruppe, Lorentz-Gruppe, $SU(n)$, Wurzelsysteme, Dynkin-Diagramme, Anwendungen auf Spektren und Teilchen. Der Band, der nach Kapitel~14 folgt, wenn man die Quarks weiterverfolgen will; die Fundstelle für Auslassung~12 in physikalischer Sprache.}

\lit{13}{Weyl, H.: \textit{Gruppentheorie und Quantenmechanik}. Hirzel 1928; 2.\ Aufl.\ 1931.}{\textbf{Auf Deutsch, das erste Buch über das Thema dieses Bandes}: unitäre Geometrie, Quantenmechanik, Gruppentheorie, dann die Anwendung --- Spin, Pauli-Prinzip, $S_n$ und $U_n$, Schur--Weyl als physikalisches Prinzip. Schwer zu lesen, aber die Quelle von allem, was danach kam.}

\lit{14}{Wigner, E.: \textit{Gruppentheorie und ihre Anwendung auf die Quantenmechanik der Atomspektren}. Vieweg 1931.}{\textbf{Auf Deutsch}: Darstellungen endlicher Gruppen, $SO(3)$, $SU(2)$, Auswahlregeln, Zeeman, Wigner--Eckart (dort erstmals), Zeitumkehr (im Anhang der englischen Ausgabe 1959). \textbf{Das Buch, das Chemiker und Physiker die Charaktertafel lesen lehrte}; die Fundstelle für Auslassung~15.}

\lit{15}{Bethe, H.: \textit{Termaufspaltung in Kristallen}. Ann.\ Phys.\ 395, 1929.}{\textbf{Auf Deutsch, die Dissertation}: Einschränkung von $SO(3)$ auf die Punktgruppen, mit den Tafeln von $O$, $T_d$, $D_4$, $D_3$ und der Einführung der Doppelgruppen. Kapitel~7 (Aufgabe~7.3), 11 und~14 im Original --- $76$ Seiten, in denen die Kristallfeldtheorie vollständig steht.}

\lit{16}{Frobenius, G.: \textit{Über Gruppencharaktere}. Sitzungsber.\ Preuß.\ Akad.\ Wiss.\ 1896, 985--1021.}{\textbf{Auf Deutsch, die Geburtsurkunde}: Charaktere als Lösungen eines Gleichungssystems aus den Klassenmultiplikationskoeffizienten, Orthogonalität, die Tafeln von $S_4$, $S_5$ und $PSL_2(\mathbb F_p)$ --- ohne Darstellungen, die erst 1897 folgen. In den \textit{Gesammelten Abhandlungen} (Springer 1968) zusammen mit den Arbeiten von 1897--1900.}

\lit{17}{Hawkins, T.: \textit{The Origins of the Theory of Group Characters}. Arch.\ Hist.\ Exact Sci.\ 7, 1971.}{Die Rekonstruktion von Frobenius' Jahr aus der Korrespondenz mit Dedekind: wie aus der Gruppendeterminante die Charaktere wurden. \textbf{Kapitel~4 (Geschichte) in der Fassung des Historikers}; Hawkins' Buch \textit{Emergence of the Theory of Lie Groups} (2000) setzt die Geschichte bis Weyl fort.}

\lit{18}{Curtis, C.\ W.: \textit{Pioneers of Representation Theory: Frobenius, Burnside, Schur, and Brauer}. AMS 1999.}{Vier Biographien und die Mathematik dazu, von einem, der sie mitgeprägt hat: Frobenius' Gruppendeterminante, Burnsides Wettlauf, Schurs Lemma, Brauers modulare Theorie. \textbf{Die Geschichtskästen dieses Bandes haben hier ihre Quelle.}}

\lit{19}{Cotton, F.\ A.: \textit{Chemical Applications of Group Theory}. 3.\ Aufl., Wiley 1990.}{\textbf{Das Lehrbuch, aus dem Generationen von Chemikern Charaktertafeln gelernt haben}: Punktgruppen, Tafeln, Molekülorbitale, Schwingungen, Kristallfeld, mit allen Tafeln im Anhang. Die Kapitel~1, 5, 6, 13 und~14 in der Sprache des Chemikers; die Beweise fehlen, die Rechnungen nicht.}

\lit{20}{Harris, D.\ C., Bertolucci, M.\ D.: \textit{Symmetry and Spectroscopy}. Oxford 1978; Dover 1989.}{Schwingungs- und Elektronenspektroskopie mit Gruppentheorie, ausführlicher als Cotton und mit den Spektren selbst: Wasser, Ammoniak, Methan, Benzol, Übergangsmetallkomplexe. \textbf{Die Kapitel~6, 9 und~13 mit den Messdaten}; die Fundstelle für die Herzberg--Teller-Analyse des Benzols.}

\lit{21}{Wilson, E.\ B., Decius, J.\ C., Cross, P.\ C.: \textit{Molecular Vibrations}. McGraw-Hill 1955; Dover 1980.}{Die Referenz für Normalschwingungen: $FG$-Methode, Symmetriekoordinaten, Auswahlregeln, Obertöne und Kombinationen, mit Beispielen durchgerechnet. \textbf{Kapitel~6 und~9 in der Fassung, die Spektroskopiker benutzen}; die Projektionsformel steht hier als Rezept.}

\lit{22}{Herzberg, G.: \textit{Molecular Spectra and Molecular Structure II: Infrared and Raman Spectra of Polyatomic Molecules}. Van Nostrand 1945.}{Alle Punktgruppen mit Tafeln und Auswahlregeln, und die Zuordnung Tausender gemessener Banden --- die Wellenzahlen dieses Bandes stammen, wo nicht anders gesagt, von hier oder aus den Nachfolgetabellen (Shimanouchi 1972). \textbf{Das Werk, das die Theorie mit dem Experiment verglich, Molekül für Molekül.}}

\lit{23}{Bishop, D.\ M.: \textit{Group Theory and Chemistry}. Clarendon 1973; Dover 1993.}{Zwischen Cotton und Serre: die Theorie mit Beweisen (Orthogonalität, Projektion, Produkte), die Anwendungen mit Rechnungen (Hückel, Schwingungen, Kristallfeld, Jahn--Teller). \textbf{Das Buch, das Chemikern die Beweise nicht erspart}; die Fundstelle für Auslassung~17.}

\lit{24}{Tinkham, M.: \textit{Group Theory and Quantum Mechanics}. McGraw-Hill 1964; Dover 2003.}{Für Physiker: Charaktere, Punktgruppen, $SO(3)$, $SU(2)$, Wigner--Eckart, Zeitumkehr, Raumgruppen und Bänder --- Kapitel~10 bis~14 in einem Band. \textbf{Die Fundstelle für Auslassung~20} und für die Darstellungstheorie der Raumgruppen, die hier nur angedeutet ist.}

\lit{25}{Hamermesh, M.: \textit{Group Theory and Its Application to Physical Problems}. Addison-Wesley 1962; Dover 1989.}{Das umfassendste der physikalischen Lehrbücher: endliche Gruppen, $S_n$ mit Young-Tableaux, kontinuierliche Gruppen, Clebsch--Gordan für $SU(n)$, Kristalle. \textbf{Für Kapitel~8 aus physikalischer Sicht} --- Hamermesh behandelt $S_n$ ausführlicher als jedes andere Physikbuch.}

\lit{26}{Sternberg, S.: \textit{Group Theory and Physics}. Cambridge 1994.}{Endliche Gruppen über Moleküle, dann $SU(2)$, $SO(3)$, Kugelflächenfunktionen, Zeitumkehr, Kristallfeld, $SU(3)$ und der Achtfache Weg, mit Anhängen zu Mathematik und Geschichte. \textbf{Der Band in anderer Reihenfolge}; die Fundstelle für Auslassungen~11 und~20.}

\lit{27}{Georgi, H.: \textit{Lie Algebras in Particle Physics}. 2.\ Aufl., Westview 1999.}{$SU(2)$, $SU(3)$, Quarks, Wurzelsysteme, $SU(5)$ --- Lie-Algebren so, wie Teilchenphysiker sie brauchen, mit Rechnungen statt Sätzen. \textbf{Kapitel~14.3 in voller Breite und die Fundstelle für Auslassung~19}; Tensormethoden für Clebsch--Gordan in $SU(3)$.}

\lit{28}{Dresselhaus, M.\ S., Dresselhaus, G., Jorio, A.: \textit{Group Theory: Application to the Physics of Condensed Matter}. Springer 2008.}{Punktgruppen, Kristallfeld, Schwingungen, Raumgruppen, Bandstruktur, Graphen und Nanoröhren --- mit Tafeln für alle $32$ Punktgruppen und Raumgruppenbeispielen. \textbf{Kapitel~14.1 in voller Länge}; die Fundstelle für Auslassung~18 und für den Dirac-Punkt als Darstellung $E''$.}
